% Part: normal-modal-logic
% Chapter: syntax-and-semantics
% Section: language-modal-logic

\documentclass[../../../include/open-logic-section]{subfiles}

\begin{document}

\olfileid{nml}{syn}{lan}

\olsection{మౌలిక మోడల్ తర్కపు భాష}

\begin{defn}
మోడల్ తర్కపు మౌలిక భాషలో ఇవి ఉంటాయి:
\begin{enumerate}
  \tagitem{prvFalse}{\tetoken{అసత్యాన్ని}{!!{falsity}} సూచించే ప్రతిజ్ఞావాక్య
    స్థిరాంకం~$\lfalse$.}{}
  \tagitem{prvTrue}{\tetoken{సత్యాన్ని}{!!{truth}} సూచించే ప్రతిజ్ఞావాక్య
    స్థిరాంకం~$\ltrue$.}{}
  \item \tetoken{లెక్కించదగిన}{!!{denumerable}s}
    \tetoken{ప్రతిజ్ఞావాక్య చరాల}{!!{propositional variable}s}
    సమితి: $\Obj p_0$, $\Obj p_1$, $\Obj p_2$, \dots
  \item ప్రతిజ్ఞావాక్య సంయోజకాలు: \startycommalist
  \iftag{prvNot}{\ycomma $\lnot$ (నిషేధం)}{}%
  \iftag{prvAnd}{\ycomma $\land$ (సంయోగం)}{}%
  \iftag{prvOr}{\ycomma $\lor$ (వికల్పం)}{}%
  \iftag{prvIf}{\ycomma $\lif$ (\tetoken{సోపాధికం}{!!{conditional}})}{}%
  \iftag{prvIff}{\ycomma $\liff$ (\tetoken{ద్విసోపాధికం}{!!{biconditional}})}{}.
  \tagitem{prvBox}{మోడల్ సంచాలకం $\Box$.}{}
  \tagitem{prvDiamond}{మోడల్ సంచాలకం $\Diamond$.}{}
\end{enumerate}
\end{defn}

\begin{defn}
మౌలిక మోడల్ భాషలోని
\emph{\tetoken{సూత్రాలను}{!!^{formula}s}}
కింది విధంగా ఆగమనాత్మకంగా నిర్వచిస్తాం:
\begin{enumerate}
\tagitem{prvFalse}{$\lfalse$ ఒక పరమాణు
  \tetoken{సూత్రం}{!!{formula}}.}{}

\tagitem{prvTrue}{$\ltrue$ ఒక పరమాణు
  \tetoken{సూత్రం}{!!{formula}}.}{}

\item ప్రతి \tetoken{ప్రతిజ్ఞావాక్య చరం}{!!{propositional variable}}
  $\Obj p_i$ ఒక (పరమాణు)
  \tetoken{సూత్రం}{!!{formula}}.

\tagitem{prvNot}{$!A$ ఒక
  \tetoken{సూత్రమైతే}{!!a{formula}}, $\lnot !A$
  కూడా \tetoken{సూత్రమే}{!!a{formula}}.}{}

\tagitem{prvAnd}{$!A$, $!B$లు
  \tetoken{సూత్రాలైతే}{!!{formula}s},
  $(!A \land !B)$ కూడా
  \tetoken{సూత్రమే}{!!a{formula}}.}{}

\tagitem{prvOr}{$!A$, $!B$లు
  \tetoken{సూత్రాలైతే}{!!{formula}s},
  $(!A \lor !B)$ కూడా
  \tetoken{సూత్రమే}{!!a{formula}}.}{}

\tagitem{prvIf}{$!A$, $!B$లు
  \tetoken{సూత్రాలైతే}{!!{formula}s},
  $(!A \lif !B)$ కూడా
  \tetoken{సూత్రమే}{!!a{formula}}.}{}

\tagitem{prvIff}{$!A$, $!B$లు
  \tetoken{సూత్రాలైతే}{!!{formula}s},
  $(!A \liff !B)$ కూడా
  \tetoken{సూత్రమే}{!!a{formula}}.}{}

\tagitem{prvBox}{$!A$ ఒక
  \tetoken{సూత్రమైతే}{!!a{formula}},
  $\Box !A$ కూడా
  \tetoken{సూత్రమే}{!!a{formula}}.}{}

\tagitem{prvDiamond}{$!A$ ఒక
  \tetoken{సూత్రమైతే}{!!a{formula}},
  $\Diamond !A$ కూడా
  \tetoken{సూత్రమే}{!!a{formula}}.}{}

\tagitem{limitClause}{ఇవి తప్ప మరేదీ
  \tetoken{సూత్రం}{!!a{formula}} కాదు.}{}
\end{enumerate}
\end{defn}

\begin{tagblock}{defNot,defOr,defAnd,defIf,defIff,defTrue,defFalse,defBox,defDiamond}
\begin{defn}
నిర్వచిత సంచాలకాలతో నిర్మించిన
సూత్రాలను కింది విధంగా అర్థం
చేసుకోవాలి:

\begin{tagenumerate}{defTrue,defFalse,defNot,defOr,defAnd,defIf,defIff,defBox,defDiamond}

\tagitem{defTrue}{$\ltrue$ అనేది
  \iftag{prvFalse}{$\lnot\lfalse$}{$(!A \lor \lnot !A)$;
    ఇక్కడ $!A$ ఒక నిర్దిష్ట పరమాణు
    \tetoken{సూత్రం}{!!{formula}}}కు సంక్షిప్త రూపం.}{}

\tagitem{defFalse}{$\lfalse$ అనేది
  \iftag{prvTrue}{$\lnot\ltrue$}{$(!A \land \lnot !A)$;
    ఇక్కడ $!A$ ఒక నిర్దిష్ట పరమాణు
    \tetoken{సూత్రం}{!!{formula}}}కు సంక్షిప్త రూపం.}{}

\tagitem{defNot}{$\lnot !A$ అనేది
  $!A \lif \lfalse$కు సంక్షిప్త రూపం.}{}

\tagitem{defOr}{$!A \lor !B$ అనేది
  \iftag{prvAnd}{$\lnot(\lnot !A \land \lnot !B)$}{$\lnot !A \lif
    !B$}కు సంక్షిప్త రూపం.}{}

\tagitem{defAnd}{$!A \land !B$ అనేది
  \iftag{prvOr}{$\lnot(\lnot !A \lor \lnot !B)$}{$\lnot (!A \lif
    \lnot !B)$}కు సంక్షిప్త రూపం.}{}

\tagitem{defIf}{$!A \lif !B$ అనేది
  \iftag{prvOr}{$\lnot !A \lor !B$}{$\lnot (!A \land \lnot !B)$}కు
  సంక్షిప్త రూపం.
  \sourcecorrection{OLTENMLLAN-001}{ఆంగ్ల మూలంలోని
  మొదటి ప్రత్యామ్నాయ వ్యక్తీకరణ చివర జతలేని
  ముగింపు కుండలీకరణ ఉంది. దాన్ని మాత్రమే
  తొలగించాం; నిషేధం, వికల్పం, రెండు
  ఉపసూత్రాలు యథాతథంగా ఉన్నాయి.}}{}

\tagitem{defIff}{$!A \liff !B$ అనేది
  $(!A \lif !B) \land (!B \lif !A)$కు
  సంక్షిప్త రూపం.}{}

\tagitem{defBox}{$\Box !A$ అనేది
  $\lnot\Diamond\lnot !A$కు సంక్షిప్త రూపం.}{}

\tagitem{defDiamond}{$\Diamond !A$ అనేది
  $\lnot\Box\lnot !A$కు సంక్షిప్త రూపం.}{}
\end{tagenumerate}
\end{defn}
\end{tagblock}

\tetoken{సూత్రం}{!!a{formula}}~$!A$లో
$\Box$ గానీ $\Diamond$ గానీ లేకపోతే,
దాన్ని \emph{మోడల్-రహిత సూత్రం} అంటాం.

\end{document}
